% ID: SEL-D0001-OA-E03
% ESTADO: validada. S002, 2026-11-07. Tres pasadas del mismo autor; ver informe.
\textbf{Datos y teoría.} Sean $C$ alojarse en la capital, $R$ alojarse en una zona rural, $H$ hospedarse en un hotel y $T$ en un apartamento turístico. Capital y zona rural forman una partición de los alojamientos, y lo mismo sucede con hotel y apartamento:
\[
P(C)=0.65,\quad P(R)=0.35,\quad P(H\mid C)=0.75,\quad P(H\mid R)=0.15.
\]
Una probabilidad condicionada se refiere sólo al grupo indicado tras la barra. La regla del producto permite pasar a la probabilidad conjunta: $P(U\cap V)=P(U)P(V\mid U)$.

\textbf{a) Probabilidad de hospedarse en un hotel.}\par\nopagebreak[4]

Aplicamos la probabilidad total, separando los casos capital y zona rural:
\[
P(H)=P(C)P(H\mid C)+P(R)P(H\mid R)
=0.65\cdot0.75+0.35\cdot0.15=\boxed{0.54}.
\]
Por tanto, la probabilidad es del $54\%$.

\textbf{b) Zona rural sabiendo que se alojó en un apartamento.}\par\nopagebreak[4]

Los datos sobre hoteles permiten calcular sus complementarios dentro de cada zona:
\[
P(T\mid C)=0.25,\qquad P(T\mid R)=0.85.
\]
La probabilidad total de apartamento es $P(T)=1-P(H)=0.46$. Por Bayes,
\[
P(R\mid T)=\frac{P(R)P(T\mid R)}{P(T)}
=\frac{0.35\cdot0.85}{0.46}
=\boxed{\frac{119}{184}\approx0.646739}.
\]
Entre quienes se han alojado en apartamentos, aproximadamente el $64.67\%$ lo ha hecho en zonas rurales. No se intercambia $P(R\mid T)$ con $P(T\mid R)$.

\textbf{Comprobación mediante una tabla conjunta.}
\[
\begin{array}{c|cc|c}
&H&T&\text{Total}\\\hline
C&0.4875&0.1625&0.65\\
R&0.0525&0.2975&0.35\\\hline
\text{Total}&0.54&0.46&1
\end{array}
\]
Todas las celdas son no negativas, suman 1 y el cociente de la celda rural-apartamento entre el total de apartamentos reproduce el resultado.
